% ============================================================
% Chapter 12 — Future Work
% ============================================================
\chapter{Future Work}
\label{ch:future}

\roadmap{This chapter turns the limitations of \cref{ch:limitations} into a
prioritized research and engineering program. It proceeds from the
scientifically most consequential extension---closing the affect-to-support
loop---through measurement upgrades, study-design improvements,
infrastructure hardening, and sensing enhancements.}

\section{Closing the loop: affect-contingent intervention}
\label{sec:future-loop}

The platform's central methodological choice---keeping sensing decoupled
from pedagogy (\cref{sec:interventions-autotrigger})---was correct for
establishing detection validity but leaves RQ2's ultimate question
unanswered: does support delivered \emph{at} moments of detected confusion
improve learning more than learner-initiated support? The natural next
study adds an affect-contingent trigger---offering the stepwise or
elaboration intervention when the windowed pipeline reports sustained
confusion---under a randomized design that compares contingent, yoked, and
learner-initiated delivery. Doing so responsibly requires the validation
of this project first, so that the trigger acts on a detector of known
fidelity, and requires a confusion-persistence threshold informed by the
transition dynamics of \citet{dmello2012dynamics}: intervene when confusion
fails to resolve, not at its onset, since early confusion can be
productive \citep{dmello2014confusion}.

\section{Human-coded ground truth}
\label{sec:future-groundtruth}

The retrospective-coding surface (\cref{sec:sensing-coding}) was built to
support, but has not yet been used for, human labeling of on-task and
affective states in the BROMP tradition \citep{ocumpaugh2015}. Recruiting
trained coders to label a sample of replays would convert this project's
convergent validation into criterion validation and yield the
inter-rater-reliable labels needed to train and benchmark
participant-independent detectors \citep{whitehill2014, bosch2015}.

\section{Study-design and analysis improvements}
\label{sec:future-study}

Three changes would strengthen the statistical program. First, add an
explicit phase/condition field to the session record so that grouped
comparisons need no external grouping (\cref{sec:methodology-stats}).
Second, formalize the pre/post design in the schema---an assessment role
field and an enforced attempt-to-assessment relation---so that gain
computation is robust rather than title-matched
(\cref{sec:methodology-gain}). Third, densify self-report, or add an
experience-sampling variant, so that the item-15 sample is larger and
better time-matched to the state windows
(\cref{sec:methodology-selfreport}).

\section{Modular learning environment and progress reporting}
\label{sec:future-modular}

Two reframed deliverables have clear completions. Wiring mastery
(\dbtable{UserMastery}) onto the prerequisite graph
(\dbtable{ProposedKC}/\dbtable{KcEdge}) would enable the runtime adaptive
learning path implied by proposal item 5 (\cref{sec:scope-item5}), and
would let the assessment surface move from adaptive feedback \emph{depth}
toward adaptive item \emph{selection}, activating the dormant BKT columns
\citep{corbett1995}. Separately, a longitudinal per-learner progress-report
generator would complete proposal item 13 from the session summaries and
mastery data already collected (\cref{sec:scope-item13}).

\section{Infrastructure hardening}
\label{sec:future-infra}

For deployment beyond a single-instance study, three fixes are prerequisite
(\cref{sec:limitations-engineering}): remediate the export-jobs
command-injection surface; move embeddings from in-memory JSON comparison
to a vector index for scale; and make the schedulers and event outbox
distributed-safe with retry and dead-lettering. Adding automatic retry to
the facial-analysis jobs would also reduce the missing-frame problem noted
in \cref{sec:sensing-caveats}.

\section{Sensing enhancements}
\label{sec:future-sensing}

Finally, the sensing layer has natural upgrades. Replacing the pupil
heuristic with remote photoplethysmography (rPPG) would add a
physiologically grounded arousal signal from the same webcam, of far
higher validity than the current proxy, following the video-heart-rate
approach of \citet{monkaresi2017}. Fitting the AOI expected weights to
task analyses rather than using placeholders would turn \allocscore{} from
a relative index into a calibrated SEEV measure
(\cref{sec:sensing-aoi}). And unifying the two affect derivations
(\cref{sec:sensing-affect}) onto the configurable server pipeline would
remove a standing source of inconsistency.

\medskip
\noindent Taken together, these extensions describe the path from a
validated measurement instrument to a closed-loop, affect-adaptive tutor---
the trajectory the conclusion draws out.
